\documentclass[10pt,a4paper]{article}
\usepackage[a4paper,margin=1.8cm]{geometry}
\usepackage[T1]{fontenc}
\usepackage{lmodern}
\usepackage{microtype}
\usepackage{booktabs}
\usepackage{tabularx}
\usepackage{array}
\usepackage[table]{xcolor}
\usepackage{enumitem}
\usepackage{titlesec}
\usepackage[colorlinks=true,linkcolor=blue!50!black,urlcolor=blue!50!black]{hyperref}

\setlist{nosep,leftmargin=1.3em}
\titlespacing*{\section}{0pt}{1.1ex}{0.6ex}
\titlespacing*{\subsection}{0pt}{0.9ex}{0.4ex}
\setlength{\parindent}{0pt}
\setlength{\parskip}{0.45ex}
\renewcommand{\arraystretch}{1.08}
\newcommand{\iss}[1]{\href{https://github.com/JieGH/RVV_code_gen_via_MLIR_xDSL/issues/#1}{\##1}}
\newcommand{\slow}[1]{\textcolor{red!70!black}{#1}}
\newcommand{\steady}[1]{\textcolor{green!45!black}{#1}}

\title{\vspace{-1.2cm}\textbf{Why B3 has two speeds, and what a smaller \texttt{nc} does}\\[2pt]
\large Brief of issues \iss{35}, \iss{36} and \iss{39}}
\author{RVV\_code\_gen\_via\_MLIR\_xDSL, \texttt{pure-xdsl} branch}
\date{27 September 2026}

\begin{document}
\maketitle
\vspace{-0.8cm}

\paragraph{Scope.} Three follow-ups from the \#31--\#34 audit. \iss{35}: how a smaller \texttt{nc} (the width of
the packed B block) changes the speed and steadiness of B3 (64$\times$384$\times$384) and the other layers with a small
M, and what B3's shipped setting should be. \iss{36}: why the vector packer is slow on B3. \iss{39}: house rules for
benchmark results. All runs on 27 September 2026 on \texttt{jbpi} (BananaPi, SpacemiT X60, VLEN=256), pinned to
\textbf{core 4}, tile 16x15. Unless marked otherwise every number is from the uninstrumented build
(\texttt{INSTRUMENTED=0}), reference protocol, \texttt{--tmin 10}. GOPs include packing. ``Mean time'' is the
reference's formula (all operations over all time); ``median'' is the median rep. Every number comes from a
committed log (list in Appendix~A).

\section{Bottom line}
\begin{enumerate}
  \item \textbf{B3 does not have a slow packer; it has a slow state.} Each B3 run starts fast and, after
  0.1--0.5\,s (vector packer) or 0.5--8\,s (scalar packer), falls into a second speed about half as fast, where it
  mostly stays. In that state \emph{both} the packing and the microkernel take longer, and the clock is unchanged:
  the core is waiting on memory. When B3 is in its fast state the vector packer is as fast as the scalar one or faster.
  \item \textbf{The vector packer's loss on B3 is how often it is in the slow state}: 66--86\% of reps at every
  \texttt{nc} from 96 to 384, against 0--46\% for the scalar packer. Both packers read B in exactly the same order;
  the vector one only sends the reads faster.
  \item \textbf{A smaller \texttt{nc} helps the scalar packer, not the shipped vector one.} Scalar: B3 is steady
  and fastest at \texttt{nc} 96--128 (7.3--7.4 GOPs, $+$11\% over 384). Vector: \texttt{nc}=128 gives 5.52 against
  4.73 at 384, still two-speed. Only \texttt{nc}$\le$64 (packed B within the 32\,KiB L1) keeps the vector packer steady.
  \item \textbf{No version of the cache model gives B3 a small \texttt{nc}}, not even with the real L2 in place of the
  missing L3. \textbf{Recommendation: drop the B3 override by default}; quote \texttt{nc}=128 only as a labelled
  hand-tuned number, and treat the slow state as the thing to fix.
  \item \textbf{House rules are written} (\iss{39}): never core 0; small gaps are compared inside one launch and
  repeated; every quoted number comes from a committed log. A new \texttt{--save-reps} option keeps every rep's time.
\end{enumerate}

\section{The slow state (\#36)}
\textbf{Every run starts fast.} The per-rep log shows the two speeds come in long stretches, not at random
(10--100 times fewer changes of speed than chance would give). First 80 reps of B3, 16x14, vector packer (F fast, s slow):
\begin{center}\small\texttt{FFFFFFFFFFFFFFFFFFFFFsFFFFFFFFFFFFFFFFFFFFFFFFFFFFFFFFFFFFFFFFFFFFFsssssssssssss}\end{center}

\textbf{In a slow rep both halves slow down.} Instrumented build (warm-up protocol, 1000 reps; debug numbers,
not comparable with the rest of this report), median per rep, fast reps against slow reps of the same run:

\begin{center}\small
\begin{tabular}{llllrr}
\toprule
tile & case & packer & fast: pack / compute (ms) & slow: pack / compute (ms) & slow reps \\
\midrule
16x14 & B3 & vector & 0.64 / 1.78 & 1.58 / 2.79 & 412 / 1000 \\
16x15 & B3 & vector & 0.75 / 1.88 & 1.46 / 2.59 & 842 / 1000 \\
16x15 & B3, \texttt{nc}=128 & vector & 0.74 / 1.58 & 1.53 / 2.12 & 533 / 1000 \\
16x16 & B3 & vector & 0.55 / 3.20 & 1.68 / 4.40 & 569 / 1000 \\
16x15 & B3, K=256 & vector & 0.33 / 1.02 & 0.83 / 1.69 & 967 / 1000 \\
16x14 & B3 & scalar & 0.76 / 1.71 & 1.24 / 1.98 & 721 / 1000 \\
16x15 & B3 & scalar & 0.69 / 1.52 & 1.11 / 1.85 & 103 / 1000 \\
16x15 & B3, \texttt{nc}=128 & scalar & 0.97 / 1.56 & --- & 0 / 1000 \\
16x15 & B3, K=256 & scalar & 0.47 / 1.00 & 1.19 / 1.57 & 981 / 1000 \\
\multicolumn{3}{l}{B2, Reuse24: both packers, all tiles} & one speed & --- & 0 / 1000 \\
\bottomrule
\end{tabular}
\end{center}
In the slow state packing takes 1.6--3.1 times as long and the kernel 16--66\% longer. Cycles per second stay at
1.63--1.81\,GHz in both states, so the core runs at full clock and stalls on memory.

\textbf{The two packers walk memory the same way.} The generated vector \texttt{pack\_CB} and the scalar one
read the same addresses of B in the same order (per 15-column panel, per k, 15 elements 1536\,B apart) and write
the same layout. The vector one uses two strided vector loads per row where the scalar one uses 15 loads, so it
has more requests in flight.

\textbf{Ruled out:} the CPU clock (governor \texttt{performance}; clock unchanged in the slow state); transparent
huge pages moving memory (\texttt{madvise} only, 0\,kB in use); B not fitting in the L2 (B3 with K=256, where B
fits, is the slowest case). \textbf{Leading guess, not tested:} a state of the shared L2 (512\,KiB, 16-way, shared by
cores 4--7) or of the hardware prefetcher that, once entered, makes this access pattern miss far more often.

\section{A smaller \texttt{nc} (\#35)}
All eleven \texttt{nc} values, the Reuse layers, B2 and the N/K variants ran inside one launch, each setting its own
dataset row; three launches per packer in the order vector, scalar, scalar, vector, vector, scalar. Average of the
three launches (lowest--highest launch in brackets). ``Slow'' is the share of reps in the slower of two clearly
separate speeds. Packed B block = 128$\times$\texttt{nc}$\times$4\,B.

\begin{center}\small
\begin{tabular}{rr|llrr|llrr}
\toprule
& & \multicolumn{4}{c|}{\textbf{vector packer (shipped)}} & \multicolumn{4}{c}{\textbf{scalar packer}} \\
\texttt{nc} & B blk & mean time & median & CV\% & slow & mean time & median & CV\% & slow \\
\midrule
32  & 16\,KiB  & \steady{6.53} (6.43--6.59) & 6.53 & 0.9  & 0\%  & 5.31 (5.21--5.48) & 5.31 & 0.6  & 0\% \\
64  & 32\,KiB  & 6.71 (6.50--7.01) & 7.55 & 16.0 & 22\% & 6.65 (6.60--6.74) & 6.65 & 1.1  & 0\% \\
96  & 48\,KiB  & 5.81 (5.48--6.34) & 5.59 & 22.4 & 66\% & \steady{7.27} (7.16--7.42) & 7.28 & 1.6  & 0\% \\
120 & 60\,KiB  & 5.67 (5.45--5.85) & 5.31 & 23.6 & 77\% & 7.26 (7.05--7.41) & 7.62 & 9.4  & 10\% \\
128 & 64\,KiB  & 5.52 (5.37--5.73) & 5.16 & 24.4 & 76\% & \steady{7.36} (7.31--7.43) & 7.58 & 7.1  & 6\% \\
192 & 96\,KiB  & 5.55 (5.44--5.69) & 4.83 & 27.4 & 76\% & 7.04 (6.68--7.65) & 7.48 & 15.0 & 37\% \\
256 & 128\,KiB & 5.10 (5.04--5.18) & 4.55 & 26.9 & 79\% & 7.07 (6.57--7.54) & 7.63 & 13.3 & 29\% \\
320 & 160\,KiB & 4.89 (4.65--5.07) & 4.34 & 28.4 & 81\% & 7.14 (6.75--7.76) & 7.82 & 13.8 & 23\% \\
384 & 192\,KiB & \slow{4.73} (4.65--4.79) & 4.27 & 30.7 & 86\% & \slow{6.65} (6.36--7.06) & 7.21 & 21.0 & 46\% \\
\bottomrule
\end{tabular}\\[2pt]
{\footnotesize B3 (64$\times$384$\times$384), GOPs. \texttt{nc} = 240 and 375 were also run and follow the same trend (issue \#35 \S6).}
\end{center}

\textbf{Other layers.} More reuse of each packed block protects against the slow state. Reuse24 (M=384) is steady
with the scalar packer at every \texttt{nc} and 0--13\% slow with the vector one; it is fastest at a large
\texttt{nc} (scalar 9.29 at 384, vector 8.53--8.56 at 240--384). Reuse8 and Reuse12 (M=128, 192) behave like B3
with the vector packer from \texttt{nc}=96 up, and are nearly steady with the scalar one. B2, the control, is
steady under both packers (CV 1--6\%) and faster at \texttt{nc}=384 (vector 11.98, scalar 11.82).

\textbf{Is it \texttt{nc} or N?} With the scalar packer and \texttt{nc} held at 128, B3 with N = 192, 256, 320,
384 gives 7.47, 7.54, 7.22, 7.36 GOPs, nearly steady (0--6\% slow). With \texttt{nc}=N the same shapes are 16--46\%
slow. So \#34's ``slow when N=384'' was \texttt{nc}=N=384. With the vector packer every variant is two-speed
(61--86\% slow).

\textbf{The cache model (host only).} \texttt{get\_optim\_mc\_nc\_kc} sizes \texttt{nc} for an L3 the board does
not have. Rerun with the real L2 (sysfs: 512\,KiB, 16-way; \texttt{bananapi.h} says 8-way):

\begin{center}\small
\begin{tabular}{llccc}
\toprule
layer & M$\times$N$\times$K & \texttt{nc} shipped (L3) & \texttt{nc} with L2, 8-way & \texttt{nc} with L2, 16-way \\
\midrule
B3 & 64$\times$384$\times$384 & 384 & 384 (model: 512) & 384 (model: 640) \\
Reuse24 & 384$\times$384$\times$384 & 384 & 128 & 192 \\
B2 & 384$\times$384$\times$64 & 384 & 256 & 256 \\
B1, B4, B5, S1--S5 & & 384--1024 & 128 & 128 \\
\bottomrule
\end{tabular}
\end{center}
The L2 version still gives B3 \texttt{nc}=384 and would instead shrink \texttt{nc} for Reuse24 and the large
layers, which the sweep shows is the wrong way for Reuse24. Swapping the L3 for the L2 is not a fix.

\section{Recommendation}
\begin{enumerate}
  \item \textbf{Drop the B3 \texttt{nc}=128 override by default.} With the shipped vector packer it is worth
  $+$17\% but leaves B3 two-speed (CV 24\%); the steadiness issue \#1 saw holds only for the scalar packer. The model
  cannot produce 128, and head-to-head comparisons need the reference's blocking. Quote an \texttt{nc}=128 number only
  as hand-tuned. (An uncommitted \texttt{config.json} edit already removes it; the five tests that expect 128 must be
  updated in the same commit as that decision.)
  \item \textbf{Chase the slow state} (\#36 next steps): per-rep L2 miss and prefetch counters in fast against slow
  reps; whether the other cores of the 4--7 cluster are active when it starts; a run on the other cluster (not
  core 0); padding \texttt{ldb} off a power of two (B3 with K=256, ties to \#38); the same on \texttt{jbpi2}.
  \item A rule-based small \texttt{nc} (``packed B within the L1'' for small-M layers with the vector packer) fits
  the data but costs the Reuse layers speed and has no known cause behind it yet. Not proposed.
\end{enumerate}

\section{House rules and tooling (\#39)}
Written in \texttt{docs/agents/execution.md}, pointed to from \texttt{README.md}:
\begin{enumerate}
  \item \textbf{Never pin to core 0} (it takes every device interrupt and runs 5--10\% slower); name the core with
  every number. Example commands and \texttt{run\_pack\_swap.sh} now use core 4. The two overnight sweep scripts keep
  core 0 only so reruns match their 19--20 September runs.
  \item \textbf{One launch is not a result.} For gaps under about 3\%, put both settings inside one launch and repeat
  over at least three launches in A-B-B-A order.
  \item \textbf{Every quoted number comes from a committed log} under \texttt{tests/pure\_xdsl\_blis\_rvv/logs/}.
  \item \textbf{Name build, protocol and core.} The default run command is now \texttt{--cpu 4}.
\end{enumerate}
\textbf{New:} \texttt{compile\_macro\_bench.sh --save-reps} writes every timed rep (with its \texttt{mc}/\texttt{nc}/\texttt{kc};
in the instrumented build, its pack and compute time) to \texttt{reps\_*.txt}, after the timed loop so nothing is added
to the timing, and gzips it on fetch. \texttt{reps\_report.py} summarises a set of launches: GOPs by mean time and by
median rep, CV, slow-rep share, how bunched the slow reps are, and how long a run stays fast. This study is the
first to follow the rules.

\textbf{Not done:} the second board, other tiles in the \texttt{nc} sweep, L2 counters, the \texttt{ldb} test.

\appendix
\section{Logs}
All under \texttt{tests/pure\_xdsl\_blis\_rvv/logs/}, \texttt{jbpi}, core 4, committed in \texttt{24c15a7}.
\begin{center}\small
\begin{tabular}{lll}
\toprule
run directory & run & config \\
\midrule
\texttt{run\_2026-09-27\_104338} & sweep, vector & \texttt{config\_issue35\_nc\_sweep.json}, \texttt{--tmin 10 --save-reps} \\
\texttt{run\_2026-09-27\_105417} & sweep, scalar & same, \texttt{--scalar-pack} \\
\texttt{run\_2026-09-27\_110456} & sweep, scalar & same, \texttt{--scalar-pack} \\
\texttt{run\_2026-09-27\_111536} & sweep, vector & same \\
\texttt{run\_2026-09-27\_112614} & sweep, vector & same \\
\texttt{run\_2026-09-27\_113653} & sweep, scalar & same, \texttt{--scalar-pack} \\
\texttt{run\_2026-09-27\_114341} & instrumented, vector & \texttt{config\_issue36\_instr\_split.json}, \texttt{--instr --save-reps} \\
\texttt{run\_2026-09-27\_115023} & instrumented, scalar & same, \texttt{--scalar-pack} \\
\bottomrule
\end{tabular}
\end{center}
Model check: \texttt{.scratch/issues/0020-model-check.c} and \texttt{-model-check.txt}. Full tables: issue files
\texttt{0020} (\#35 \S6) and \texttt{0021} (\#36, Results).
\end{document}
